\section{Fluctuations of the satisfiability transition}\label{sec:introduction}

Add independent random clauses to a Boolean formula until the formula
becomes unsatisfiable. How much does the number of clauses at this first
failure vary from one sample to another? For every fixed clause size
$k\ge4$, we prove that its variance has order $n$, where $n$ is the
number of variables. For three-literal clauses the variance lies between
positive constant multiples of $n$ and $n\log n$.

A \emph{proper $k$-clause} is a disjunction of literals on $k$ distinct
variables. Fix $k\ge3$ and $n\ge k$, and let $C_1,C_2,\ldots$ be
independent uniform choices from the $2^k\binom nk$ proper signed
clauses on $n$ variables. Signs are independent and fair; whole clauses
may repeat. Put
\[
 F_m=C_1\wedge\cdots\wedge C_m,\qquad F_0=\mathrm{true},\qquad
 H_n=\min\{m\ge1:F_m\text{ is unsatisfiable}\}.
\]
For a fixed $B>0$, the capped first failure time is
$T_{n,B}=\min\{H_n,\lfloor Bn\rfloor\}$.
Each assignment survives a clause with probability $1-2^{-k}$, so
\begin{equation}\label{eq:intro-first-moment}
 \Pr(H_n>m)\le2^n(1-2^{-k})^m.
\end{equation}
Thus $H_n$ is almost surely finite and has finite moments.

\begin{theorem}\label{thm:variance}
Fix an integer $k\ge3$, and write
\[
 \ell_k(n)=\begin{cases}\log(en),&k=3,\\1,&k\ge4,\end{cases}
 \qquad U_k=\frac{\log2}{-\log(1-2^{-k})}.
\]
There are constants $C_k,c_k>0$ such that
\[
 \Var(H_n)\le C_kn\ell_k(n)\quad(n\ge k),\qquad
 \Var(H_n)\ge c_kn\quad\text{for all sufficiently large }n.
\]
For every fixed $B>0$ there is $C_{k,B}<\infty$ such that
\[
 \Var(T_{n,B})\le C_{k,B}n\ell_k(n)\quad(n\ge k).
\]
If $B>U_k$, then also $\Var(T_{n,B})\ge c_kn$ for all sufficiently
large $n$, with the size cutoff allowed to depend on $k$ and $B$.
\end{theorem}

All constants concern fixed $k$. The lower condition on the cap is
substantive: an early cap can stop the process before its random
transition. In the proof, it is enough to have integer caps
$L_n=O(n)$ with $\Pr(H_n>L_n-1)\le1/4$ eventually;
$B>U_k$ is an explicit sufficient condition.
This argument leaves a logarithmic gap for three-literal clauses; the
companion paper~\cite{LinearThreeSATVariance} closes it and proves
$\Var(H_n)=\Theta(n)$ for the uncapped proper 3-clause process.

The identity $\Pr(F_m\text{ is satisfiable})=\Pr(H_n>m)$ turns these
variance bounds into bounds on the transition window. Between any fixed
probability levels $\eta$ and $1-\eta$, $0<\eta<1/2$, its width is
$O_{k,\eta}(\sqrt{n\ell_k(n)})$. Chebyshev's inequality also gives
\[
 \Pr\bigl(|H_n-\EE H_n|\ge t\sqrt{n\ell_k(n)}\bigr)
 \le C_kt^{-2}\qquad(t>0).
\]
This is concentration about the finite-size expectation. The limiting
threshold $\alpha_k$ of \cite[Theorem~\ThresholdTheorem]{FixedKThreshold} identifies
$\lim_n\EE H_n/n$, but supplies no bound on the bias
$\EE H_n-\alpha_kn$ at the fluctuation scale above.
For $k=3$, the separate companion~\cite[Theorem~1.1]{ComputableThreeSATThreshold}
proves that the limiting threshold $\alpha_3$ is a computable real.

\subsection{Earlier work and the new estimate}

The location and the width of a threshold are different questions.
Friedgut~\cite[Theorem~1.3]{Friedgut} proved a sharp threshold sequence
for each fixed clause size, without asserting convergence of that
sequence. Achlioptas and Peres~\cite[Theorems~1--2]{AchlioptasPeres}
used a weighted second moment to locate the threshold within $O(k)$ of
$2^k\log2$ as $k$ grows. Ding, Sly, and Sun~\cite[Theorem~1]{DSS}
proved the predicted limiting threshold for all sufficiently large $k$.
Carenini's paper~\cite[Corollary~1.2]{Carenini2026Polynomial}, made public
on October 5, 2026, establishes the satisfiability conjecture for every
fixed $k\ge3$. We credit Carenini with priority for the resolution of the
satisfiability conjecture.
The companion \cite[Theorem~\ThresholdTheorem]{FixedKThreshold} establishes a limiting
location for every fixed $k\ge3$; its independent forcing argument gives
the variance bound $O_k(n^{1+2/k})$ for the capped
last-satisfiable prefix $\min(H_n-1,2^{k+1}n)$
\cite[Proposition~\ThresholdVariance]{FixedKThreshold}.

For transition width, Wilson~\cite[Corollary~4]{Wilson} proved an
$\Omega_k(\sqrt n)$ separation of central quantiles in the independent
proper-clause model. We convert that separation into the variance lower
bound, keeping the integer endpoints so the argument allows atoms.
Carenini~\cite[Corollaries~7.10--7.11]{Carenini} proves an
$O_{k,\eta}(n/\log n)$ window upper bound for fixed $k\ge3$, including
proper independent clauses. Her subsequent paper gives the polynomial
window bound $O_{k,\eta}(n^{1/2+1/k})$ and the hitting-time variance bound
$O_k(n^{1+2/k})$~\cite[Theorems~1.1 and~1.3]{Carenini2026Polynomial};
the window estimate, together with the Abbe--Montanari criterion, yields
the limiting-threshold theorem credited above.
The upper bounds here sharpen these fluctuation estimates. For $k=3$,
the companion~\cite{LinearThreeSATVariance} further removes the logarithmic
loss in this paper's upper bound.

Our upper bound uses the coordinate-omission form of the Efron--Stein
inequality~\cite{EfronStein,BBLM}. Deleting one clause leaves its $k$
variables available for root-flip tests of the remaining solutions.
When no other clause meets two of these variables, the $k$ tests are
independent and use clauses of length $k-1$; we account separately for
such collisions. A set of assignments that is easily
killed by those shorter clauses must also have a short expected lifetime
under $k$-clauses. The replacement of a shorter constraint by a block of
native clauses has a qualitative predecessor in Friedgut's half-cube
argument~\cite[Lemma~5.7]{Friedgut}; sequential conditioning on a residual
assignment set is developed in Carenini~\cite[Theorem~4.3 and Lemma~6.1]{Carenini}.
We prove a block estimate uniform in the number of variables and the
block length. It also covers assignment sets that no single $k$-clause
can kill.

The main improvement comes after this replacement step. Instead of
bounding the number of clauses incident to the deleted clause's variables
by a logarithmic cutoff, we retain that number as a random parameter.
Its fixed moments are bounded. The proof must therefore use the
incidence count without destroying the independence of the shorter-clause
tests or of future ordinary clauses. Conditioning on the incidence
pattern through the cap achieves both. Collisions involving two of the
deleted clause's variables are paid for by joint incidence and collision
moments.

\subsection{The route through the upper bound}

Write $L=\lfloor Bn\rfloor$. If deleting $C_i$ while preserving the
other clause indices delays the capped first failure by $D_i$, then
\[
 \Var(T_{n,B})\le\sum_{i=1}^{L}\EE D_i^2.
\]
The delay $D_i$ counts the prefixes at which restoring $C_i$ destroys
satisfiability, so $D_i^2$ counts ordered pairs of such prefixes. The
problem is to control this occupation time.

Here is the shorter-clause estimate used for that purpose. A collection
of clauses \emph{kills} $S\subseteq\{0,1\}^u$ if no member of $S$
satisfies every clause. Let $q_d(S)$ be the probability that $d$ independent
proper $(k-1)$-clauses kill $S$, and let $\tau_k(S)$ be the expected
number of independent proper $k$-clauses through its first kill. For
$u\ge k$, $d\ge1$, and $\beta=k/(k-1)$, we prove in
Section~\ref{sec:replacement} that
\[
 \tau_k(S)q_d(S)^\beta\le C_kd^\beta.
\]
This estimate holds for every assignment set, with no lower cutoff on
$q_d(S)$.

Section~\ref{sec:deletion} connects this statement to clause deletion.
Call the variables of $C_i$ its roots, and fix them to the unique values
that falsify $C_i$. Let $S_m$ be the nonroot parts of solutions of the
prefix with $C_i$ omitted and the roots fixed this way. A clause disjoint
from the roots is an ordinary $k$-clause on the nonroots. A clause meeting
one root and satisfied by that root hides a $(k-1)$-clause: its nonroot
part becomes a constraint when that root is flipped. If restoring $C_i$
destroys satisfiability, each allowed root flip must fail. In the absence
of collisions this gives $k$ independent tests killing $S_m$.

Conditional on the exposed incidence pattern, let $Z$ be the number of
clauses meeting the roots and put $d=\max\{1,Z\}$. The lifetime estimate
bounds both the expected remaining duration of a deletion event and the
total time spent at levels $q_d(S_m)\ge a$. Write $q=q_d(S_m)$.
The $k$ independent tests contribute $q^k$, while the expected remaining
duration costs a factor $q^{-\beta}$. Integrating the resulting
$q^{k-\beta}$ over time costs a second lifetime factor: the expected
occupation above level $a$ is bounded by a constant times $a^{-\beta}$.
The near-zero part of the resulting integral is therefore
\[
 \int a^{\,k-2\beta-1}\,da.
\]
At $k=3$, $k=2\beta$ and the finite cap produces a logarithm; for
$k\ge4$, $k>2\beta$ and the integral is bounded. Section~\ref{sec:occupation}
then averages the root counts and proves the capped upper bound.
Exponential tails remove a sufficiently late cap. Wilson's quantile
bound supplies the separate lower-bound argument.

Section~\ref{sec:interfaces} removes a sufficiently late cap and records
the exact first-unsatisfiable and last-satisfiable endpoint identities.
Section~\ref{sec:lower-bounds} proves the lower bound, completing the
variance theorem. Section~\ref{sec:three-clause-refinements} specializes
the cap constants to three clauses and transfers the proper-clause
concentration to the Poisson clause law used in the threshold companion.
After that main argument,
Appendix~\ref{sec:replacement-sharpness} determines the limits of the
two estimates for arbitrary assignment sets.
Appendix~\ref{sec:direct-comparison} compares proper-clause models at
different sizes, and Appendix~\ref{sec:appendix-c} gives a second
three-clause proof based on prefix exposure and incidence truncation.
